\documentclass[11pt]{article}
\usepackage[T1]{fontenc}
\usepackage[utf8]{inputenc}
\usepackage{lmodern}
\usepackage[a4paper,margin=27mm]{geometry}
\usepackage{amsmath,amssymb,amsthm,booktabs,array}
\usepackage{microtype}
\usepackage[hidelinks]{hyperref}
\usepackage{url}
% Standard article section numbering.


\hypersetup{pdftitle={King on generalized staircases with equal odd blocks},
            pdfauthor={Xiaojian Zeng}}
\newtheorem{theorem}{Theorem}
\newtheorem{lemma}[theorem]{Lemma}
\newtheorem{proposition}[theorem]{Proposition}
\theoremstyle{remark}
\newtheorem{remark}[theorem]{Remark}
\DeclareMathOperator{\mex}{mex}
\newcommand{\N}{\mathbb{Z}_{\geq0}}
\newcommand{\parity}[1]{\overline{#1}}
\title{King on generalized staircases with equal odd blocks}
\author{Xiaojian Zeng\\[2pt]
\small PBC School of Finance, Tsinghua University, 100190\\
\small\texttt{xiaojianzeng@foxmail.com}}
\date{}
\begin{document}
\maketitle
\begin{abstract}
Let $D_k=((kr)^r,((k-1)r)^r,\ldots,r^r)$, where exponents denote
repeated parts and $r$ is positive and odd. In the impartial King game,
a move deletes the first row, the first column, or both, and must leave
a nonempty partition. We prove Conjecture 5.1 of the January 2025 arXiv
version of Gottlieb, Krnc and Mur\v si\v c: writing $q=k\bmod(r+2)$,
the Sprague--Grundy value of $D_k$ is $0$ for odd $q$, $1$ for positive
even $q$, and $2$ for $q=0$. The proof determines an entire boundary
word by a square transfer. An explicit array with three parity regions
satisfies the local mex recurrence, including the exceptional corner
and both alternating phase transitions. The argument is uniform in $r$
and $k$ and has no computational premise. The later SSRN text has not
been compared; no priority claim is made.
\end{abstract}

\section{The game and the assertion}

For a nonempty partition $\lambda=(\lambda_1\geq\cdots\geq\lambda_\ell>0)$,
let $\lambda[i,j]$ be the tuple
$(\lambda_{i+1}-j,\ldots,\lambda_\ell-j)$ with all nonpositive entries
removed. Following Gottlieb, Krnc and Mur\v si\v c
\cite[Definitions 2.1 and 3.1]{gkm2025}, the King moves are
\[
 K=\{(0,1),(1,0),(1,1)\},
\]
and a move is legal only if its destination is nonempty. We use normal
play: the player with no legal move loses. In the Sprague--Grundy
framework, the recursive value is the least nonnegative integer absent
from the successor values; see Sprague \cite[Section 6, Satz I]{sprague1935}.
Thus
\begin{equation}\label{eq:game}
 G(\lambda)=\mex\{G(\lambda[i,j]):(i,j)\in K,
                                      \ \lambda[i,j]\ne\varnothing\}.
\end{equation}
In particular $\mex\varnothing=0$. Every move decreases the number of
cells, so the recursion is well defined, and $0\leq G(\lambda)\leq3$.
The empty partition is neither a position nor an option.

For positive integers $r,k$, write
\[
 D_k=((kr)^r,((k-1)r)^r,\ldots,r^r),
\]
where each exponent indicates repetition of a part. Throughout the
proof $r$ is fixed and odd.

\begin{theorem}\label{thm:main}
For every positive odd integer $r$ and every positive integer $k$, let
$q$ be the residue of $k$ modulo $r+2$, with $0\leq q<r+2$. Then
\begin{equation}\label{eq:main}
 G(D_k)=
 \begin{cases}
 0,&q\text{ is odd},\\
 1,&q\text{ is even and nonzero},\\
 2,&q=0.
 \end{cases}
\end{equation}
\end{theorem}

This is exactly \cite[Conjecture 5.1]{gkm2025}. The same source gives
the rectangle formula in Theorem 3.8, attributed there to Berlekamp;
Lemma 3.13 gives $G(D_k)=(k-1)\bmod3$ when $r=1$ and $k\geq1$;
and Lemma 3.14 gives value $2$ at $k=1$ and $1$ at $k\geq2$ for
equal even blocks. The contribution here is a uniform proof of the
odd-block assertion, with an explicit boundary and square-interior
description. These comparisons refer to the inspected arXiv version
\cite{gkm2025}. A later SSRN posting \cite{gkmssrn2025} was identified,
but its full text has not been compared. Accordingly, the result is
identified by its precise source statement, without a claim of priority
or of the conjecture's current status elsewhere.

\section{Boundary data and square transfer}

Set
\begin{equation}\label{eq:parameters}
\begin{gathered}
 p=r+2,\qquad t=1+((k-1)\bmod p),\qquad
 h=\left\lfloor\frac{k-1}{p}\right\rfloor\bmod2,\\
 L=p-t,\qquad x=L\bmod2=(t+1)\bmod2,\qquad y=1-x,\\
 a=3-h,\qquad b=2+h.
\end{gathered}
\end{equation}
The parameters $t,h,L,x,y,a,b$ depend on $k$. For $0\leq j<r$, define
\begin{equation}\label{eq:word}
 W_k(j)=
 \begin{cases}
 2,&j=0\text{ and }L=0,\\
 x,&j<L\text{ and }j\text{ is even},\\
 a,&j<L\text{ and }j\text{ is odd},\\
 b,&j\geq L,\ j\text{ is even},\ (j,L)\ne(0,0),\\
 y,&j\geq L\text{ and }j\text{ is odd}.
 \end{cases}
\end{equation}
We will prove the stronger assertion
\begin{equation}\label{eq:boundary}
 B_k(j):=G(D_k[0,j])=W_k(j)\qquad(0\leq j<r).
\end{equation}
The extra phase $h$ is needed for the boundary: for $r=3$ the formula
gives $B_1=(0,3,0)$ and $B_6=(0,2,0)$, although the root values agree.

Transposing a Young diagram exchanges the row and column moves and
preserves the diagonal move. Induction on the number of cells therefore
shows that conjugate partitions have equal $G$-values. Since $D_k$ is
self-conjugate, write $A_k(i,j)=G(D_k[i,j])$ whenever the indicated
partition is nonempty; then $A_k(i,j)=A_k(j,i)$.

Direct deletion of complete blocks gives
\begin{equation}\label{eq:shapes}
\begin{aligned}
 D_k[r,0]=D_k[0,r]&=D_{k-1}&& (k\geq2),\\
 D_k[r,r]&=D_{k-2}&& (k\geq3).
\end{aligned}
\end{equation}
These are shape identities, not additional legal moves. Deletions compose,
so the first $r\times r$ square has prescribed edges
\begin{equation}\label{eq:edges}
\begin{aligned}
 A_k(r,s)=A_k(s,r)&=B_{k-1}(s)&& (0\leq s<r),\\
 A_k(r,r)&=B_{k-2}(0)&& (k\geq3).
\end{aligned}
\end{equation}
For $k=2$ the corner $(r,r)$ is absent. Every cell with $0\leq i,j<r$
exists. Its options are its right, down and diagonal neighbors that
exist. Descending induction on $i+j$ uniquely determines all interior
values from \eqref{eq:edges}. It suffices, therefore, to exhibit values
with these edges and the correct mex at every interior cell.

\section{A local mex array}

The following construction is algebraic; it does not define an infinite
game. For any $L\in\N$, put $x=L\bmod2$ and $y=1-x$, and choose
$(a,b)$ to be either $(3,2)$ or $(2,3)$. For $i,j\in\N$ let
$m=\min(i,j)$, $n=\max(i,j)$. Define $F(i,j)$ by the following tables,
with rows indexed by the parity of $m$ and columns by the parity of $n$,
in the order even, odd:
\begin{equation}\label{eq:array}
\begin{array}{c@{\qquad}c@{\qquad}c}
 n<L&m\leq L\leq n&L<m\\[2pt]
 P=\begin{pmatrix}x&a\\b&y\end{pmatrix}&
 C=\begin{pmatrix}b&y\\x&a\end{pmatrix}&
 T=\begin{pmatrix}y&b\\a&x\end{pmatrix}.
\end{array}
\end{equation}
At the single point $(L,L)$ override the table by $F(L,L)=2$.

\begin{lemma}\label{lem:array}
For all $i,j,L\in\N$ and both choices of $(a,b)$,
\begin{equation}\label{eq:local}
 F(i,j)=\mex\{F(i+1,j),F(i,j+1),F(i+1,j+1)\}.
\end{equation}
\end{lemma}
\begin{proof}
By symmetry suppose $i\leq j$. Successor triples below are ordered
as (down, right, diagonal); repeated entries do not affect mex.
If $i<j$ and all four cells use $P$, the cases are
\begin{center}
\begin{tabular}{cccc}
\toprule
$i$ parity&$j$ parity&$F(i,j)$&Successor triple\\
\midrule
even&even&$x$&$(b,a,y)$\\
even&odd&$a$&$(y,x,b)$\\
odd&even&$b$&$(x,y,a)$\\
odd&odd&$y$&$(a,b,x)$\\
\bottomrule
\end{tabular}
\end{center}
Here $\{x,y\}=\{0,1\}$ and $\{a,b\}=\{2,3\}$, so each row
has the required mex. On an ordinary $P$ diagonal the value and triple
are $x$ and $(a,a,y)$ for even $i$, or $y$ and $(b,b,x)$ for odd $i$.
The $T$ cases follow by exchanging $x$ with $y$ and $a$ with $b$.
If all four cells use $C$, they contain all four values, so the
successor values have mex equal to the current value.

The remaining off-diagonal cases are listed in Table~\ref{tab:seams}.
Each row is used only when its coordinate conditions are feasible.
The two exceptional diagonal cases are
\[
\begin{array}{c|c|c}
 \text{location}&F(i,j)&\text{successor triple}\\\hline
 i=j=L-1\quad(L\geq1)&1&(0,0,2)\\
 i=j=L&2&(1,1,0).
\end{array}
\]
Again every displayed triple has the asserted mex for both choices
of $(a,b)$.

To check exhaustiveness, for $i<j$ the range $j\leq L-2$ is ordinary
$P$, and $j=L-1$ is the first group in Table~\ref{tab:seams}. For
$j=L$, the range $i\leq L-2$ is ordinary $C$, leaving $(L-1,L)$.
For $j>L$, the three cases $i<L$, $i=L$, $i>L$ give ordinary $C$,
the last group of the table, and ordinary $T$, respectively. On the
diagonal only $L-1$ and $L$ are exceptional. This also covers $L=0,1$
by omitting infeasible cases.
\end{proof}

\begin{table}[tb]
\centering
\caption{The off-diagonal cases meeting a region boundary in
\eqref{eq:array}.}\label{tab:seams}
\begin{tabular}{llccc}
\toprule
Location&$L$ parity&Other parity&$F(i,j)$&Successor triple\\
\midrule
$j=L-1,\ i<j$&even&$i$ even&$a$&$(1,b,0)$\\
&even&$i$ odd&$1$&$(a,0,b)$\\
&odd&$i$ even&$1$&$(b,0,a)$\\
&odd&$i$ odd&$b$&$(1,a,0)$\\
\midrule
$i=L-1,\ j=L$&even&---&$0$&$(2,a,1)$\\
&odd&---&$0$&$(2,b,1)$\\
\midrule
$i=L,\ j>L$&even&$j$ odd&$1$&$(0,b,a)$\\
&even&$j$ even&$b$&$(a,1,0)$\\
&odd&$j$ even&$1$&$(0,a,b)$\\
&odd&$j$ odd&$a$&$(b,1,0)$\\
\bottomrule
\end{tabular}
\end{table}

\section{Closing the transfer}

\subsection*{The rectangle base}
For positive side lengths $u,v$, the King value of a rectangle is
\begin{equation}\label{eq:rectangle}
 R(u,v)=
 \begin{cases}
 0,&u,v\text{ both odd},\\
 2,&u,v\text{ both even},\\
 1,&u,v\text{ of opposite parity and }\min(u,v)\text{ odd},\\
 3,&u,v\text{ of opposite parity and }\min(u,v)\text{ even}.
 \end{cases}
\end{equation}
This is \cite[Theorem 3.8]{gkm2025}; a short verification makes the
base self-contained. A unit side gives a path with alternating values
$0,1$, starting with $R(1,1)=0$. For $u,v\geq2$, use induction on
$u+v$ and symmetry to take $u\geq v$. If $u>v$, the successor sets
are, in the parity order odd--odd, even--even, odd--even, even--odd,
\[
 \{1,2,3\},\quad\{0,1,3\},\quad\{0,1,2\},\quad\{0,2,3\}.
\]
For $u=v$ odd or even they are $\{2,3\}$ or $\{0,1\}$,
respectively. Their mex values give \eqref{eq:rectangle}.

Now $D_1[0,j]$ is an $r$-by-$(r-j)$ rectangle. Its value is $0$
for even $j$ and $3$ for odd $j$. Since $k=1$ gives
$L=r+1$, $x=0$, $a=3$, this proves \eqref{eq:boundary} at $k=1$.
We do not assert the $F$-interior formula at this base case.

\subsection*{Ordinary phases and the absent corner}
Assume $k\geq2$ and \eqref{eq:boundary} for all smaller positive
indices. Form $F$ using the current parameters \eqref{eq:parameters}.
First suppose $2\leq t\leq p$. The preceding index has parameter
$t-1$, the same $h$, threshold $L+1$, and interchanged low values $x,y$.
Since $r$ is odd, substitution into \eqref{eq:array} gives
\begin{equation}\label{eq:ordinaryedge}
 F(r,s)=F(s,r)=
 \begin{cases}
 y,&s\leq L\text{ and }s\text{ even},\\
 a,&s\leq L\text{ and }s\text{ odd},\\
 b,&s>L\text{ and }s\text{ even},\\
 x,&s>L\text{ and }s\text{ odd},
 \end{cases}\qquad(0\leq s<r).
\end{equation}
These are precisely $W_{k-1}(s)$.

If $t\geq3$, then $L<r$ and the odd--odd entry of $T$ gives
$F(r,r)=x$. The parameter of $k-2$ is $t-2$, so
$W_{k-2}(0)=(t-1)\bmod2=x$. If $t=2$, then $L=r$ and the
override gives $F(r,r)=2$; when $k>2$, the parameter of $k-2$ is $p$
and its root is also $2$. Thus in these cases the entire prescribed
boundary agrees with $F$. Lemma~\ref{lem:array} and backward uniqueness
give the actual values throughout the first square.

At $k=2$, the corner $(r,r)$ is absent. Only $(r-1,r-1)$ has it
as an option. The two remaining options both have value
$W_1(r-1)=0$, since $r-1$ is even, and the proposed interior value
is $1$. Removing the fictitious corner value $2$ changes the option
set from $\{0,2\}$ to $\{0\}$; both have mex $1$. Every other
interior identity is unchanged. This handles $k=2$, including $r=1$,
without assigning a value to the empty partition.

\subsection*{Both phase transitions}
It remains to consider $t=1$ and $k>1$. Write $k=\ell p+1$,
$\ell\geq1$. Here $L=r+1$ is even, $x=0$, $y=1$, and
$h=\ell\bmod2$. Every point of $[0,r]^2$ uses $P$.
The preceding boundary has parameter $p$ and phase $1-h$; for $s>0$
its values are $a$ at even $s$ and $1$ at odd $s$. These agree with
$F(r,s)$. At $s=0$, however, the preceding root is $2$, whereas
$F(r,0)=F(0,r)=a$. The corner satisfies $F(r,r)=1$, as required:
the parameter of $k-2$ is the positive even integer $p-1$.

If $h=1$, then $a=2$, and all prescribed boundary entries already
agree with $F$. If $h=0$, then $a=3$, $b=2$. Replace the two
boundary entries $(r,0),(0,r)$ by $2$. For $r\geq3$ the only affected
interior cells are $(r-1,0),(0,r-1)$. At the latter cell the proposed
value is $0$ and the successor triple after replacement is $(2,2,1)$,
whose mex is still $0$. Symmetry treats the other cell. For $r=1$,
these interior cells coincide at $(0,0)$ and the same triple applies.
Thus the corrected boundary preserves every interior recurrence in
both phases.

We have proved, for every $k\geq2$,
\begin{equation}\label{eq:interior}
 G(D_k[i,j])=F(i,j)\qquad(0\leq i,j<r).
\end{equation}
At $i=0$, \eqref{eq:array} gives exactly \eqref{eq:word}, including
the override when $L=0$. Together with the rectangle base, this proves
\eqref{eq:boundary} by induction on $k$.

\begin{proof}[Proof of Theorem~\ref{thm:main}]
Take $j=0$ in \eqref{eq:boundary}. If $t=p$, the value is $2$.
Otherwise it is $x=(t+1)\bmod2$. The parameter $t$ represents
$k\bmod p$, with its zero residue represented by $p$. These are
exactly the three cases of \eqref{eq:main}.
\end{proof}

\section{Exact checks and reproducibility}

The proof above is symbolic. Exact computation was used to detect
transcription errors and to check the treatment of phase changes and
missing destinations. The accompanying verification files contain a
block-transfer implementation, a direct evaluation on entire Young
diagrams, and a separately implemented recursive evaluator that forms
the partition tuples after each legal deletion. No floating-point
arithmetic is used for game values.

The independent direct implementation checks all
$1\leq k\leq3(r+2)+2$ for $r\in\{1,3,5,7,9\}$, and
\[
 k\in\{1,2,p-1,p,p+1,2p,2p+1\},\qquad p=r+2,
\]
for $r\in\{15,37,39\}$. Its saved output records 136 boards,
27,464,908 evaluated cells, and 26,800 comparisons in the first squares,
with no mismatch. The tuple evaluator checks 245 subpartitions, namely
all first-square entries for $r\in\{1,3,5\}$ and $1\leq k\leq7$.
These are exact finite observations, with precisely the stated coverage.

The local identities were also encoded as 14 counterexample queries
over unbounded integers in the SMT solver Z3 \cite{demoura2008}; each
returned \texttt{unsat}. They test the mex array, edges, corners,
phase transitions, rectangle recurrence, and extraction of the stated
formula. This checks the encoded identities; the interpretation of
partitions and the complete induction remain the written proof above.
No formal proof-assistant certification of the full theorem is claimed.

Exact predicates and separately implemented verifiers also appear in
the finite-library certification of Zeng et al.\ \cite{zeng2026finite}.
There exhaustive computations establish assertions within specified
ternary libraries. In the present problem, the finite game evaluations
corroborate a proof valid for all positive odd $r$ and all positive $k$.
This distinction determines what each computation can certify.

The supplied reproduction guide identifies the self-contained scripts,
saved integer-query files and their execution commands. The checks used
Python 3.10.12 and, for the symbolic queries, Z3 4.15.3. The proof does
not address unequal block dimensions or other chess movesets.

\paragraph{Computational and AI assistance.}
OpenAI Codex assisted with the boundary description, proof development,
verification code, source searches, and preparation of this manuscript.
A separate model context reconstructed the argument and implemented
additional checks. That separation is a review procedure, not statistical
independence or external peer review. The mathematical argument is given
in full, and the exact computations and their limits are disclosed above.

\bibliographystyle{plain}
\bibliography{references}
\end{document}
